\chapter{Magnetic Flux Tubes in a Dense Lattice}
\label{ch:periodic}
\acresetall

\begin{summary}
If the superconducting nuclear material of a neutron star contains magnetic flux tubes, 
the magnetic field is likely to vary rapidly on the scales where \acs{QED} effects 
are important.
In this chapter, I construct a cylindrically symmetric toy model of a flux 
tube lattice in which the influence of neighbouring flux tubes is taken into account.
We compute the effective action densities using the \acs{WLN} technique.
The numerics predict a greater effective energy density in the 
region of the flux tube, but a smaller energy density in the regions 
between the flux tubes compared to a locally-constant-field approximation. 
We also compute the interaction energy between a flux tube and its neighbours 
as the lattice spacing is reduced from infinity. Because our flux tubes 
exhibit compact support, this energy is entirely non-local and
predicted to be zero in local approximations
such as the derivative expansion. This quantity can take 
positive or negative values depending on the magnetic field profile and the 
specific definition of the interaction energy. 
\end{summary}

\section{Introduction}


Computing the quantum effective action for magnetic flux tube 
configurations is a problem that has generated considerable interest
and has been explored through a variety of approaches~\cite{Gornicki1990271, 
1998MPLA...13..379S, 0264-9381-12-5-013, PhysRevD.51.810, 1999PhRvD..60j5019B,
PhysRevD.62.085024, 2001PhRvD..64j5011P, 2003PhRvD..68f5026B, 
2005NuPhB.707..233G, 2006JPhA...39.6799W, Weigel:2010pf, Langfeld:2002vy}
(see section \ref{sec:EAisoflux}). 
Partly, this 
is because it is a relatively simple problem for analyzing non-homogeneous 
generalizations of the Heisenberg-Euler action and for exploring 
limitations of techniques such as the derivative expansion. But, this 
is also a physically important problem because tubes of magnetic flux 
are very important for the quantum mechanics of electrons due 
to the Aharonov-Bohm effect, and they appear 
in a variety of interesting physical scenarios such as stellar astrophysics, 
cosmic strings, in superconductor vortices, and quark confinement
\cite{2003PrPNP..51....1G}.

In the present context, we are concerned with the role that magnetic 
flux tubes play in the superconducting nuclear material of compact stars.
In this scenario, the \ac{QED} effects are particularly interesting because 
the magnetic flux tubes, if they exist, are confined to tubes which may be 
only a few percent of the Compton wavelength, $\lambda_{\rm C}$, in radius. Specifically, 
the flux must be confined to within the London penetration depth of the superconducting 
material, which for neutron stars has been estimated to be 80 fm $=$
0.032 $\lambda_e$~\cite{PhysRevLett.91.101101}. Moreover, the 
flux tube density is expected to be proportional to the 
average magnetic field. For a background field near the quantum critical 
strength, $B_k$, such as in a neutron star, the distance between flux tubes 
is comparable to a Compton wavelength. This Compton wavelength scale is 
also the scale at which the non-locality of \ac{QED} becomes important and 
at which powerful local techniques like the derivative expansion 
are no longer appropriate for computing the effective action.

The free energy associated with these flux tubes is a factor in determining 
whether the nuclear material of a neutron star is a type-I or type-II 
superconductor. The free energy of a flux tube is determined by 
looking at the energies associated with the magnetic field, with the 
creation of a non-superconducting region in the superconductor, and 
with interactions between the flux tubes. Flux tubes can only form 
if it is energetically favourable to do so compared to expelling the 
field due to the Meissner effect. In this chapter, I would 
like to explore the contribution from \ac{QED} to this free 
energy. For an isolated flux tube, this is an additional source 
of energy for creating the magnetic field. For a lattice of 
flux tubes, there is also an energy contribution from the 
presence of neighbouring flux tubes because of the non-local 
nature of quantum field theory.

The energy of two flux tubes has been previously
computed using \ac{WLN} methods and for flux tubes with aligned fields, the 
energy is larger than twice the energy of a single flux tube when the flux tubes 
are closely spaced~\cite{Langfeld:2002vy}. This result implies that there is a 
repulsive interaction between the flux tubes due to \ac{QED} effects,
strengthening the likelihood of the type-II scenario in neutron stars. This interaction 
energy increases as the flux tubes are placed closer together, and can have a similar 
magnitude as the \ac{QED} correction to the energy when the flux tubes are closely spaced.


In this chapter, I will further explore the nature of this 
phenomenon in \ac{QED} using
 a highly parallel implementation of the \ac{WLN} 
technique implemented on a \ac{GPU} architecture. Specifically, I explore 
cylindrically symmetric magnetic field profiles for isolated flux 
tubes and periodic profiles designed to model properties of a 
triangular lattice. This 
algorithm cannot be applied to spinor \ac{QED} calculations 
in our model lattice because 
of the well-known fermion problem of \ac{WLN}. However, the problem 
does not affect \ac{ScQED} calculations. Therefore, 
we explore the quantum-corrected energies of 
isolated flux tubes for both scalar and spinor 
electrons and use this comparison to speculate 
about the relationship of our cylindrical lattice model 
and the spinor \ac{QED} energies of an Abrikosov lattice 
of flux tubes that may be found in neutron stars.





\subsection{Cylindrical Magnetic Fields} 

We consider our flux tubes to have cylindrical symmetry so that 
the field points along the \bhattext{z}-direction with a 
profile that depends on the radial coordinate, $\rho$. 
In this case, we have $\vec{B} = B(\rho)$\bhattext{z}. 
We can find the magnetic field in terms of the vector potential 
from $\vec{B} = \curl{\vec{A}}$ and the gauge choice that $A_0 =
A_\rho = A_z = 0$. Then,

\be 
	\label{eqn:Bperiodic} B(\rho) = \frac{A_\phi(\rho)}{\rho} +
	\frac{dA_\phi(\rho)}{d\rho}. 
\ee 

The vector potential, $\Aphi$, can be characterized by a profile function, 
$f_\lambda$, and a dimensionless flux parameter, $\mathcal{F}=\frac{e}{2\pi}F$:

\be 
	\Aphi = \frac{\mathcal{F}}{e \rho}f_\lambda(\rho).
\ee 
If we choose $f_\lambda(\rho=0) = 0$, then the total flux within
in a radius of $L_\rho$ is 
\be 
	\Phi=\frac{2\pi}{e}\mathcal{F} f_\lambda(L_\rho).  
\ee 

In terms of the profile function, the magnetic field is written
\be
	\label{eqn:Bzfromfl}
	B_z(\rho)=\frac{\mathcal{F}}{e\rho}\frac{df_\lambda(\rho)}{d\rho} 
		=\frac{2\mathcal{F}}{e}\frac{df_\lambda(\rho^2)}{d(\rho^2)}. 
\ee 


\section{Isolated Flux Tubes}

To explore isolated magnetic flux tubes, we consider the following 
profile function:

\be
	f_\lambda(\rho^2) = \frac{\rho^2}{\rho^2 + \lambda^2}.
\ee
From equation (\ref{eqn:Bzfromfl}), this give a magnetic field 
with a profile
\be
	\label{eqn:isoB}
	B_z(\rho^2) = \frac{2\mathcal{F}}{e}\frac{\lambda^2}{(\rho^2+\lambda^2)^2}.
\ee
This profile is a smooth flux tube representation that can be evaluated quickly.
Moreover, flux tubes 
with this profile were studied previously using \ac{WLN}
\cite{Moyaerts:2003ts, MoyaertsLaurent:2004}.



\section{Cylindrical Model of a Flux Tube Lattice}

\begin{figure}
	\centering
		\includegraphics[width=10cm]{images/latticesymmetry.eps}
	\caption[Cylindrical model of a hexagonal lattice]
	{In a type-II superconductor, there are neighbouring flux 
	tubes arranged in a hexagonal Abrikosov lattice which have a nonlocal 
	impact on the effective action of the central flux tube (left). In our model, 
	we account for the contributions from these neighbouring flux tubes in 
	cylindrical symmetry by including surrounding rings of flux (right).}
	\label{fig:latticesymmetry}
\end{figure}

In a neutron star, we do not have isolated flux tubes. The tubes are likely arranged in a dense lattice with the 
spacing between tubes on the order of the Compton wavelength, with the size of a flux tube a few percent of 
the Compton wavelength. Specifically, the maximum size of a flux tube is on the order of the 
coherence length of the superconductor, which for neutron stars has been estimated to be 
$\xi \approx 30$ fm~\cite{PhysRevLett.91.101101}. 
This situation can be directly computed in the \ac{WLN} technique. However, 
this requires us to integrate over two spatial dimensions instead of one. Moreover, it requires the use of more 
loops to more precisely probe the spatial configurations of the magnetic field. Despite these problems, 
it is very interesting to consider a dense flux tube lattice. Unlike the isolated flux tube, the wide-tube limit 
of the configuration doesn't have zero field, but an average, uniform background field. If this background field is 
the size of the critical field, there are interesting quantum effects even in the wide-tube limit. 

In this section, I build a cylindrically symmetric toy model of a hexagonal flux tube lattice. We focus 
on one central flux tube and treat the surrounding six flux tubes as a continuous ring with six units 
of flux at a distance $a$ from the central tube. The next ring will contain twelve units of flux at a 
distance of $2a$, etc (see figure \ref{fig:latticesymmetry}). 
Because of this condition, the average strength of the field is fixed, and the field becomes uniform in the 
wide tube limit instead of going to zero. For small values of $\lambda$, we will have nonlocal 
contributions from the surrounding rings in addition to the local contributions from the central 
flux tube.

It is difficult to construct a model of this scenario if the flux tubes bleed into one another 
as they are placed close together. For example, with Gaussian flux tubes or flux tubes 
with the profile used in the previous section, it is difficult to increase the width of the flux tubes 
while accounting for the magnetic flux that bleeds out of their regions. Moreover, it is 
difficult to integrate these schemes to find the profile function $f_\lambda(\rho)$ which is needed 
to compute the scalar part of the Wilson loop.  In order to keep each tube as a distinct entity which 
stays within its assigned region, we assign 
a smooth function with compact support to represent each tube. 
This is most easily done with the bump function, $\Psi(x)$, defined as
\be
	\Psi(x) = 
	\begin{cases}
	e^{-1/(1-x^2)} & \mbox{ for } |x| < 1\\
	0 & \mbox{ otherwise} 
	\end{cases}.
\ee
This function can be viewed as a rescaled Gaussian.
%Unfortunately, by the identity theorem, smooth functions with compact support cannot be 
%analytic unless they vanish everywhere.


We start by defining the magnetic field outside of the central flux tube. Here, the magnetic field 
is a constant background field, with the flux ring contributing a bump of width $\lambda$. The height 
of the bump must go to zero as the width of the flux tube approaches the distance between flux tubes, 
and should become infinite as the flux tube width goes to zero:

\be
	B_z(\rho>\frac{a}{2}) = B_{\rm bg} + A\left(\frac{a-\lambda}{\lambda}\right)\left[\Psi(2(\rho-n a)/\lambda)-B\right].
\ee
with $n\equiv \lfloor\frac{\rho+a/2}{a}\rfloor$.


If we require 6 units of flux in the first outer ring, 12 in the second, 
and so on (see figure \ref{fig:latticesymmetry}), the size of the background field 
is fixed to $B_{\rm bg} = \frac{6 \mathcal{F}}{ea^2}$. The total flux contribution due to the $\lambda$-dependent 
terms must be zero:

\be
	\int_{(n-1/2)a}^{(n+1/2)a} \rho A\left(\frac{a-\lambda}{\lambda}\right)
		\left[\Psi(2(\rho-n a)/\lambda)-B\right] d\rho = 0
\ee

\be
	\frac{\lambda}{2}\int_{-1}^1 \left(\frac{\lambda}{2}x+na\right)\Psi(x)dx-Ba^2n = 0.
\ee
This fixes the value of the constant $B$ to
\be
	B=\frac{q_1}{2}\frac{\lambda}{a}.
\ee
The numerical constant $q_1$ is defined by
\be
	q_1 = \int_{-1}^{1}\Psi(x)dx \approx 0.443991.
\ee

For a given bump amplitude, $A$, the magnetic field will become negative if $\lambda$ becomes small enough.
Therefore, we replace $A$ with its maximum value for which the field is positive if $\lambda > \lambda_{\rm min}$
for some choice of minimum flux tube size:

\be
	A=\frac{12 \mathcal{F}}{e a q_1 (a-\lambda_{\rm min})}.
\ee
The choice of $\lambda_{\rm min}$ sets the tube width at which the field between the flux tubes vanishes. If 
$\lambda < \lambda_{\rm min}$, the magnetic field between the flux tubes will 
point in the $-\hat{\vec{z}}$-direction.

Because we are trying to fit a hexagonal peg into a round hole, we must treat the central flux tube differently. 
For example, the average field inside the central region for a unit of flux, is different than the average 
field in the exterior region. Therefore, even when $\lambda \rightarrow a$, the field cannot be quite uniform. 
We consider the field in the central region to be a constant field with a bump centered at $\rho = 0$:

\be
	B_z(\rho < \frac{a}{2}) = A_0 \Psi(2\rho/\lambda) + B_0.
\ee
The constant term, $B_0$ is fixed by requiring continuity with the exterior field at $\rho = a/2$:

\be
	B_0 = \frac{6\mathcal{F}}{e a^2}\left(1-\frac{a-\lambda}{a-\lambda_{\rm min}}\right).
\ee
The bump amplitude, $A_0$, is fixed by fixing the flux in the central region,
\be
	\int_0^{a/2}\rho\left[A_0\Psi(2\rho/\lambda) + B_0\right]d\rho = \frac{\mathcal{F}}{e}:
\ee

\be
	A_0 \left(\frac{\lambda}{2}\right)^2\int_0^1 x \Psi(x)dx + 
		\frac{B_0}{2}\left(\frac{a}{2}\right)^2 = \frac{\mathcal{F}}{e}
\ee

\be
	A_0 = \frac{4 \mathcal{F}}{\lambda^2 e q_2}\left(1-\frac{3}{4}
		\left(1-\frac{a-\lambda}{a-\lambda_{\rm min}}\right)\right),
\ee
where the numerical constant, $q_2$, is defined by

\be
	\label{eqn:q2}
	q_2 \equiv \int_0^1 x \Psi(x) dx \approx 0.0742478.
\ee

Finally, collecting together the important expressions, 
the cylindrically symmetric flux tube lattice model is

\ba
\label{eqn:ffbless}
B_z(\rho\le\frac{a}{2}) &=& \frac{4 \mathcal{F}}{\lambda^2 e q_2}
	\left(1-\frac{3}{4}\left(\frac{\lambda-\lambda_{\rm min}}{a-\lambda_{\rm min}}\right)\right)\Psi(2\rho/\lambda) \nonumber \\
	& & + \frac{6\mathcal{F}}{ea^2}\left(\frac{\lambda-\lambda_{\rm min}}{a-\lambda_{\rm min}}\right)
\ea

\be
\label{eqn:Bzext}
B_z(\rho>\frac{a}{2}) = \frac{6 \mathcal{F}}{e a^2}\left(\frac{\lambda-\lambda_{\rm min}}{a-\lambda_{\rm min}}\right)
	+ \frac{12 \mathcal{F}}{q_1ea\lambda}\left(\frac{a-\lambda}{a-\lambda_{\min}}\right)
	\Psi(2(\rho-na)/\lambda)  
\ee
with
\be
	\Psi(x) = 
	\begin{cases}
	e^{-1/(1-x^2)} & \mbox{ for } |x| < 1\\
	0 & \mbox{ otherwise} 
	\end{cases},
\ee

\be
	n=\left \lfloor\frac{\rho+a/2}{a}\right\rfloor,
\ee

\be
	\label{eqn:q1}
	q_1 \equiv \int_{-1}^{1}\Psi(x)dx \approx 0.443991,
\ee
and
\be
	q_2 \equiv \int_0^1 x \Psi(x) dx \approx 0.0742478.
\ee
The magnetic field profile defined by these equations is shown in figures
\ref{fig:periodicB} and \ref{fig:B3d}. The current density required to 
created fields with this profile is shown in figure \ref{fig:current}.

\begin{figure}
	\centering
		\includegraphics[width=10cm]{images/periodicB.eps}
	\caption[Magnetic field in a cylindrical lattice model]
	{The cylindrical lattice flux tube model for several 
	values of the width parameter $\lambda$. Here we have taken $a =\sqrt{8}\lambda_e$
	and $\lambda_{\rm min} = 0.1a$. Note that in the limit $\lambda\rightarrow a$
	the field is nearly uniform with a mound in the central region. This is a 
	consequence of the flux conditions in cylindrical symmetry requiring different
	fields in the internal and external regions. The height of the $\lambda=0.1a$
	flux tube extends beyond the height of the graph to about $61.9B_k$. 
	A 3D surface plot of the $\lambda=0.6a$ field profile is shown in figure 
	\ref{fig:B3d}.}
	\label{fig:periodicB}
\end{figure}

\begin{figure}
	\centering
		\includegraphics[width=12cm]{images/B3d.eps}
	\caption[Surface plot of magnetic field model]
	{A 3D surface plot of the $\lambda = 0.6a$ (red) magnetic field profile from 
	figure \ref{fig:periodicB}.}
	\label{fig:B3d}
\end{figure}

\begin{figure}
	\centering
		\includegraphics[width=12cm]{images/current.eps}
	\caption[Classical current required to create magnetic field model]
	{The current densities required to create the $\lambda=0.6a$ (red) magnetic field profile. 
	The current is given by the curl of the magnetic field, $J_\psi(\rho) = -\frac{d B_z(\rho)}{d\rho}$.
	The conversion to SI units is $1B_k/\lambda_e \approx 6\times 10^{38}{\rm A/m}^3$.}
	\label{fig:current}
\end{figure}

\subsection{The Classical Action}

The classical action, $\Gamma^0$, is infinite for this configuration. To obtain 
a finite result, we must look at the action per unit length in the z-direction, 
per unit time, and per flux tube region (\ie for $\rho < a/2$). 
The action for such a region in 
cylindrical coordinates is given by 

\be
	\frac{\Gamma^0}{T L_z} = -\pi \int_0^{a/2}\rho B_z(\rho)^2d\rho.
\ee
We substitute in equation (\ref{eqn:ffbless}), the magnetic field in the interior 
region:
\ba
	\frac{\Gamma^0}{T L_z} &=& -\pi \int_0^{a/2}\rho \biggr[
	\frac{4\mathcal{F}}{\lambda^2e^2q_2}\left(1-\frac{3}{4}
	\left(\frac{\lambda-\lambda_{\rm min}}{a-\lambda_{\rm min}}\right)\right)
	\Psi(2\rho/\lambda)\nonumber \\
	& & + \frac{6\mathcal{F}}{ea^2}\left(\frac{\lambda-\lambda_{\rm min}}{a-\lambda_{\rm min}}\right)
	\biggr]^2 d\rho.
\ea

After some algebra, we are left with an expression for the classical action,

\ba
	\frac{\Gamma_0}{TL_z} &=& \pi \int_0^{a/2}\rho B_z(\rho)^2d\rho \\
	& = & -\frac{\pi \mathcal{F}^2}{e^2a^2}\biggr\{4\frac{a^2q_3}{\lambda^2q_2^2}
	+\left(\frac{\lambda-\lambda_{\rm min}}{a - \lambda_{\rm lmin}}\right)
	\biggr[ \nonumber \\
	& &\left(\frac{9}{4}\frac{a^2q_3}{\lambda^2q_2^2}-\frac{9}{2}\right)
	\left(\frac{\lambda-\lambda_{\rm min}}{a - \lambda_{\rm lmin}}\right) 
	-6\frac{a^2q_3}{\lambda^2q_2^2}+12 \biggr]\biggr\},
\ea
with $q_3$ being another numerical constant related to integrating the bump function:

\be
	\label{eqn:q3}
	q_3 \equiv \int_0^1 x\left(\Psi(x)\right)^2 dx 
	= \int_0^1 x e^{-\frac{2}{1-x^2}}dx \approx 0.0187671.
\ee

\subsection{Integrating to Find the Potential Function}

To compute the Wilson loops, it is generally required to use the 
vector potential which describes the magnetic field. For us, this means that 
we must find $f_\lambda(\rho)$ for our magnetic field model. 
This could always be done numerically, but can be computationally costly since 
it is evaluated by every discrete point of every worldline in the ensemble.
For computations on the \ac{CUDA} device, an increase in the complexity of the 
kernel often means that less memory 
resources are available per processing thread, limiting 
the number of threads that can be computed concurrently.
It is therefore preferable to find an analytic expression for this function.
From equation (\ref{eqn:Bzfromfl}), this function is related to the integral 
of the magnetic field with respect to $\rho^2$. For the inner region, we have

\ba
	f_\lambda(\rho < a/2) &=& \frac{e}{2\mathcal{F}}\biggr[\frac{4\mathcal{F}}{\lambda^2eq_2}
		\left(1-\frac{3}{4}\left(\frac{\lambda-\lambda_{\rm min}}{a-\lambda_{\rm min}}\right)\right)
		\int_0^{\rho^2}\Psi\left(\frac{2\rho'}{\lambda}\right)d \rho'^2 \nonumber \\
		& &+\frac{6 \mathcal{F}}{ea^2}\left(\frac{\lambda-\lambda_{\rm min}}{a-\lambda_{\rm min}}\right)\rho^2\biggr].
\ea
The integral over the bump function can be computed in terms of the exponential integral 
$E_i(x) = \int_{-\infty}^x \frac{e^t}{t} dt$:
\ba
	\int_0^{\rho^2}\Psi\left(\frac{2\rho'}{\lambda}\right)d\rho'^2 = 
	\begin{cases}
	\left(\frac{\lambda}{2}\right)^2\biggr[2q_2+\left(\frac{4\rho^2}{\lambda^2}-1\right) 
		e^{-\frac{1}{1-\frac{4\rho^2}{\lambda^2}}} & {}\\
		~~~~~~~~~~- E_i\left(-\frac{1}{1-\frac{4\rho^2}{\lambda^2}}\right)\biggr]& \mbox{ for } \rho < \lambda/2\\
	\frac{q_2\lambda}{2} & \mbox{ for } \rho \ge \lambda/2 
	\end{cases}.
\ea
So, our expression for the profile function in the inner region is
\be
	f_\lambda(\rho\le a/2) = \left(1-\frac{3}{4}\left(\frac{\lambda-\lambda_{\rm min}}{a-\lambda_{\rm min}}\right)\right)
		\Phi(2\rho/\lambda) + \frac{3\rho^2}{a^2}\left(\frac{\lambda-\lambda_{\rm min}}{a-\lambda_{\rm min}}\right),
\ee
with
\be
	\Phi(x) \equiv
	\begin{cases}
	1 + \frac{1}{2q_2}\left(x^2-1\right)e^{-\frac{1}{1-x^2}}
		-\frac{1}{2q_2}E_i\left(-\frac{1}{1-x^2}\right)& \mbox{ for } x < 1\\
	1 & \mbox{ for } x \ge 1 
	\end{cases}.
\ee

The exterior integral is a bit more challenging, but we can make significant progress 
and obtain an approximate expression. The first term is a constant given by the value of the profile 
function at $\rho = a/2$. This value is given by the flux in the central flux tube, which 
we have already chosen to be 1,
\be
	f_\lambda(\rho>a/2) = 1 + \frac{e}{\mathcal{F}}\int_{a/2}^{\rho}\rho'B(\rho'>a/2)d\rho'.
\ee
We may plut the magnetic field, equation (\ref{eqn:Bzext}), into this expression to get
\ba
	f_\lambda(\rho>a/2) &=& 1 + \frac{3}{4}\left(\frac{4\rho^2}{a^2}-1\right)
		\left(\frac{\lambda-\lambda_{\rm min}}{a-\lambda_{\rm min}}\right) \nonumber \\
		& &+\frac{12}{q_1a\lambda}\left(\frac{a-\lambda}{a-\lambda_{\rm min}}\right)
		\int_{a/2}^{\rho}\rho'\Psi\left(\frac{2(\rho-na)}{\lambda}\right)d\rho'.
\ea
The remaining integral is over every bump between $\rho'=a/2$ and $\rho'=\rho$. We express the result 
as a term which accounts for each completely integrated bump, and an integral over the partial bump if $\rho$ is 
within a bump:
\ba
	f_\lambda(\rho>a/2) &=& 1 + \frac{3}{4}\left(\frac{4\rho^2}{a^2}-1\right)
		\left(\frac{\lambda-\lambda_{\rm min}}{a-\lambda_{\rm min}}\right)
		+3n(n-1)\left(\frac{a-\lambda}{a-\lambda_{\rm min}}\right) \nonumber \\
		& &+\frac{3\lambda}{q_1a}\left(\frac{a-\lambda}{a-\lambda_{\rm min}}\right)\chi(2(\rho-na)/\lambda),
\ea
where
\be
	\chi(x_0) = 
	\begin{cases}
	0 & \mbox{ for } x_0 \le -1  \\
	\int_{-1}^{x_0}xe^{-\frac{1}{1-x^2}}dx 
			+ \frac{2na}{\lambda}\int_{-1}^{x_0}e^{-\frac{1}{1-x^2}}dx& \mbox{ for } |x_0| < 1\\
	 \frac{2naq_1}{\lambda} & \mbox{ for } x_0 \ge 1 	  
	\end{cases}.
\ee
One of the integrals in $\chi(x_0)$ can be expressed in terms of the exponential integral:
\be
	\int_{-1}^{x_0}xe^{-\frac{1}{1-x^2}}dx =\frac{1}{2}\left[(x_0^2-1)e^{-\frac{1}{1-x_0^2}}
		-E_i\left(-\frac{1}{1-x_0^2}\right)\right].
\ee

The remaining integral can't be simplified analytically. To use this integral in our numerical model, 
it must be computed for each discrete point on each loop for each $\rho_{\rm cm}$ and $T$ value. 
Therefore, it is worthwhile to consider an approximate expression which models the integral, and 
can be computed faster than performing a numerical integral each time. To find this approximation, 
I computed the numerical result at 300 values between $x_0=-1.2$ and $x_0=1.2$. The data was then 
input into Eureqa Formulize, a symbolic regression program which uses genetic algorithms to find 
analytic representations of arbitrary data~\cite{formulize}. 
A similar technique has been used to produce 
approximate analytic solutions of ODEs~\cite{geneticODEs}.
The result is a model of the numerical data points with a maximum error of 0.0001 on the range $|x_0| < 1.0$:
%\be
%	\int_{-1}^{x_0} e^{-\frac{1}{1-x^2}}dx \approx 0.218 + \frac{0.393x_0\cosh{(0.806x_0^4-0.696x_0^2+0.0902)}}
%		{\cosh{(0.825x_0^2-0.0234x_0+0.375)}}
%\ee

\be
	\int_{-1}^{x_0} e^{-\frac{1}{1-x^2}}dx \approx 
	\frac{0.444}{1+\exp{\left[ -3.31x_0 - 
	\frac{5.25x_0^3 - 3.31x_0^2\sin{(x_0)}\cos{(-0.907x_0^2-1.29x_0^8)}}{\cos{(x_0)}}\right]}}.
\ee

This function evaluates a factor of ten faster than the numerical integral evaluated at the 
same level of precision with the \ac{GSL} Gaussian quadrature library functions and 
with fewer memory registers. 
Using this approximation
introduces a systematic uncertainty  which is small compared to that associated with the 
discretization of the loop integrals, and considerably smaller than the statistical error bars.

Using these expressions for the integrals, we can express $f_\lambda(\rho)$ in any region in terms of 
exponential integrals, exponential, and trigonometric functions with suitable precision:

\small
\ba
	\chi(x_0) \approx 
	\begin{cases}
	0 & \mbox{ for } x_0 \le -1 \\
	\frac{1}{2}\left[(x_0^2-1)e^{-\frac{1}{1-x_0^2}}
		-E_i\left(-\frac{1}{1-x_0^2}\right)\right] +
		\frac{2na}{\lambda}\biggr[\\
		~~~~~\frac{0.444}{1+\exp{\left( -3.31x_0 - 
	\frac{5.25x_0^3 - 3.31x_0^2\sin{(x_0)}\cos{(-0.907x_0^2-1.29x_0^8)}}
		{\cos{(x_0)}}\right)}}\biggr] & \mbox{ for } |x_0| < 1\\
	 \frac{2naq_1}{\lambda} & \mbox{ for } x_0 \ge 1 	  
	\end{cases}.
\ea
\normalsize
For computation, we may express the exponential integral 
as a continued fraction:

\be
E_i(-x) = -E_1(x) = -e^x\left[ \cfrac{1}{x+1-\cfrac{1}{x+3-\cfrac{4}{x+5-\cfrac{9}{x+7-...}}}}\right].
\ee
This expression converges very rapidly in the range of 
interest (i.e. $x>1$) using Lentz's 
algorithm~\cite{nr3rd}.

\begin{figure}
	\centering
		\includegraphics[width=10cm]{images/periodicfl.eps}
	\caption[The profile function in a cylindrical lattice model]
	{The function $f_\lambda(\rho)$ for the above described magnetic 
	field model. The flux conditions require the function to pass through the 
	black dots. A quadratic function corresponds to a uniform field while 
	a staircase function corresponds to delta-function flux tubes. The parameter 
	$\lambda$ smoothly transitions between these two extremes.
	Each of these functions corresponds to a magnetic field profile in 
	figure \ref{fig:periodicB}.}
	\label{fig:periodicfl}
\end{figure}

The profile function, $f_\lambda(\rho)$, is plotted in figure \ref{fig:periodicfl}.


\section{Results}
\label{sec:periodic_results}

\subsection{Comparing Scalar and Fermionic Effective Actions}

Because of the fermion problem of \ac{WLN} 
(see section \ref{sec:fermionproblem}), the 1-loop effective 
action for the cylindrical flux tube lattice model could not be 
computed for the case of spinor \ac{QED}. 
The fermion problem does not 
affect the scalar case. So, we are confined to analyzing this 
model for \ac{ScQED}. 

For isolated flux tubes, the decay of the magnetic field for large distances
protects the calculations from the fermion problem (see section 
\ref{sec:fermionproblem}). Therefore, the
effective action can be computed for both scalar and spinor 
\ac{QED}. In figure \ref{fig:scalfermiratio}, we plot the 
ratio of the spinor to scalar 1-loop correction term for 
identical magnetic fields, along with the prediction of the 
\ac{LCF} approximation for large values of $\lambda$. The 
\ac{LCF} approximation in \ac{ScQED} is given by
\ba
	\label{eqn:LCFscal}
	\Gamma^{(1)}_{\rm scal} &=& -\frac{1}{2\pi}\int_0^\infty dT 
	\int_0^\infty \rho_{\rm cm} d\rho_{\rm cm}\frac{e^{-m^2T}}{T^3} \nonumber \\
	& &\left\{\frac{eB(\rho_{\rm cm})T}{\sinh{(eB(\rho_{\rm cm})T)} }
	- 1 +\frac{1}{6}(eB(\rho_{\rm cm})T)^2\right\}.
\ea
This can be compared to the spinor \ac{QED} approximation, 
equation (\ref{eqn:LCFferm}).


\begin{figure}
	\centering
		\includegraphics[width=10cm]{images/scalfermiratio.eps}
	\caption[Comparison of 1-loop term in \acs{QED} and \acs{ScQED}]
	{The ratio of the 1-loop term in \ac{QED} to the 1-loop term 
	in \ac{ScQED}. The solid line is the \ac{LCF} approximation, 
	while the points are the result of \ac{WLN} calculations. Note that the 
	\ac{LCF} approximation breaks down near $\lambda = \lambda_e$ and 
	that the statistics from point to point are strongly correlated. This plot shows that 
	the 1-loop correction in \ac{ScQED} differs from the \ac{QED} correction by a factor 
	close to unity for a wide range of flux tube widths.}
	\label{fig:scalfermiratio}
\end{figure}

There are two important notes to make about figure \ref{fig:scalfermiratio}. 
Firstly, the \ac{LCF} approximation is only a good approximation 
for $\lambda \gg \lambda_e$, and isn't accurate when pushed near its formal 
validity limits~\cite{MoyaertsLaurent:2004}. The second note is that the 
statistics of the \ac{WLN} points are strongly correlated. So, we conclude 
that the \ac{ScQED} 1-loop correction is larger than the \ac{QED} correction 
for large $\lambda$, and this appears to be reversed for small $\lambda$. 
However, the large \ac{WLN} error bars and the invalidity of the 
\ac{LCF} approximation near $\lambda = \lambda_e$ prevent us from 
seeing how this transition happens. Nevertheless, the main 
conclusion from this figure is that the scalar 1-loop correction 
reflects the behaviour of the full \ac{QED} 1-loop correction 
to within a factor of about 2 over a wide range of flux tube widths 
for isolated flux tubes.

Besides using a finite field profile, the fermion problem can also be circumvented by increasing
the electron mass. 
From equation (\ref{eqn:cylEA}), we can see that the square of the electron mass 
sets the scale for the exponential suppression of the large proper time Wilson loops that 
contribute to the fermion problem. However, if we increase the fermion mass, we are 
reducing the Compton wavelength of our theory so that the flux tube lattice is no longer 
dense in terms of the (new) Compton wavelength. It is the Compton wavelength of the theory that 
determines what is meant by `dense'. We therefore cannot avoid the fermion 
problem for dense lattice models by changing the electron mass.

Based on the results presented in figure \ref{fig:scalfermiratio}, we conclude that the coupling between the electron's spin 
and the magnetic field do not have a dramatic effect on the vacuum 
energy for isolated flux tubes. Therefore, we expect that \ac{ScQED} provides a good model of the 
underlying vacuum physics near these flux tubes, at least at the level our toy model 
flux tube lattice. 


\subsection{Flux Tube Lattice}

The \ac{WLN} technique computes an effective action density which is then integrated to 
obtain the effective action. This quantity differs from the Lagrangian in that it is 
not determined by local operators, but encodes information about the field everywhere
through the worldline loops.
Like the classical action, the 1-loop term of the effective action per unit length 
is infinite for a flux tube lattice because the field extends infinitely far. 
For this reason, we define the effective action to be the action density integrated 
over the region of a central flux tube $(0<\rho<a/2)$:

\small
\ba
	\frac{\Gamma}{\mathcal{T}L_z} &=& -\pi \int_0^{a/2}\rho B_z(\rho)^2d\rho \nonumber \\
	& &-\frac{1}{2\pi}\int_0^{a/2}\rho_{\rm cm} d\rho_{\rm cm}\int_0^\infty \frac{dT}{T^3}e^{-m^2T}
	\nonumber \\
	& &\times \left\{ \mean{W}_{\rho_{\rm cm}} - 1 +\frac{1}{6}(eB(\rho_{\rm cm})T)^2\right\}.
\ea
\normalsize

The 1-loop 
term of the effective action density is plotted in figure \ref{fig:EAscffpeek} for the 
cylindrical flux tube lattice model. The most pronounced feature of this density is 
that there is a negative contribution from the regions where the field is strong. 
This contribution has the same sign as the classical term. Therefore, quantum correction 
tends to reinforce the classical action. A less pronounced feature is that there is a 
positive contribution arising from the $\rho > \lambda/2$ region, in between 
the lumps of magnetic field which represent the flux tubes. 
In this region, the local magnetic field is positive, but small. 




\begin{figure}
	\centering
		\includegraphics[width=10cm]{images/EAscffpeek.eps}
	\caption[The 1-loop action density for a cylindrical lattice in \acs{ScQED}.]
	{The \ac{ScQED} effective action density for the central 
	flux tube in our cylindrical lattice model for several tube widths, 
	$\lambda$. The average field strength is the critical field, $B_K$. 
	The effective action is positive in between flux tubes due to nonlocal 
	effects.}
	\label{fig:EAscffpeek}
\end{figure}

To interpret this feature, we consider the relative contributions between the 
Wilson loops and the counter term. These terms are shown in figure 
\ref{fig:CounterTermDomination} for the constant field case. 
For all values of proper time, $T$, the counter terms dominate, giving 
an overall negative sign. In order for the action density to be positive, there 
must be a greater contribution from the Wilson loop average than from the 
counter term, since this term tends to give a positive contribution to the 
action. In our flux tube model, this seems to occur in the regions between the 
flux tubes. In these regions, the local contribution from the counter term is 
relatively small because the field is small. However, the contribution from the 
Wilson loop average is large because the loop cloud is exploring the nearby 
regions where the field is much larger. The effect is largest where the field 
is small, but becomes large in a nearby region. We therefore interpret the positive 
contributions to the 1-loop correction from these regions as a non-local 
effect. A similar example of such an effect from fields which vary on 
scales of the Compton wavelength has been observed previously using the 
\ac{WLN} technique~\cite{PhysRevD.84.065035}. 


\begin{figure}
	\centering
		\includegraphics[width=10cm]{images/CounterTermDominationScalar.eps}
	\caption[Wilson loop and counter term contributions for a constant field]
	{The Wilson loop and counter term contributions to the integrand 
	of the effective action for a constant field in \ac{ScQED}. For constant fields, the 
	effective action is always negative due to the domination of the counter 
	term over the Wilson loop. For non-homogeneous fields, 
	A positive effective action density signifies 
	that non-local (\ie $T > 0$) effects dominate the counter term.}
	\label{fig:CounterTermDomination}
\end{figure}


\begin{figure}
	\centering
		\includegraphics[width=10cm]{images/EAscffvsl.eps}
	\caption[The 1-loop \acs{ScQED} term versus flux tube width]
	{The 1-loop \ac{ScQED} term of the effective action as a function 
	of flux tube width, $\lambda/a$, for several values of the flux tube spacing, $a$.
	The dotted lines are computed from the \ac{LCF} approximation.}
	\label{fig:EAscffvsl}
\end{figure}

In figure \ref{fig:EAscffvsl}, we plot the magnitude of the 1-loop 
\ac{ScQED} term of the effective 
action as a function of the flux tube width. As the flux tubes become smaller, there 
is an amplification of the 1-loop term, just as there is for the classical action. 
Similarly, for more closely spaced flux tubes, $a$ is smaller, and the 1-loop term 
increases in magnitude. The ratio of the 1-loop term to the classical term is plotted 
in figure \ref{fig:ActRatscff}. The quantum contribution is greatest for closely spaced, 
narrow flux tubes, but does not appear to become a significant fraction of the total action.

\begin{figure}
	\centering
		\includegraphics[width=10cm]{images/EAscffresid.eps}
	\caption[The deviation of the cylindrical flux tube lattice from the \ac{LCF}
	approximation]
	{The residuals between the \ac{WLN} results and the \ac{LCF} approximation 
	for the cylindrical flux tube lattice. The level of agreement is not expected because 
	the field varies rapidly on the Compton wavelength scale. This agreement is believed to 
	be due to an averaging effect of integrating over the electron degrees of freedom making a
	mean-field approximation appropriate.}
	\label{fig:EAscffresid}
\end{figure}

We observe that the \ac{LCF} approximation is surprisingly good despite the 
fact that the magnetic field is varying rapidly on the Compton wavelength scale of the electron.
We plot the residuals showing the deviations between the \ac{WLN} results and the \ac{LCF}
approximation in figure \ref{fig:EAscffresid}. To understand this, recall the discussion 
surrounding figure \ref{fig:CounterTermDomination}. The Wilson loop term is sensitive to the 
average magnetic field through the loop ensemble, $\mean{B}_{\rm e}$. In contrast, the counter 
term is sensitive to the magnetic field at the center of mass of the loop, $B_{\rm cm}$. 
Since these terms carry opposite signs, we can understand the difference from the 
constant field approximation in terms of a competition between these terms. When 
$B_{\rm cm} < \mean{B}_{\rm e}$, such as when the center of mass is in a local minimum of the 
field, there is a reduction of the energy relative to the locally constant field case, with 
a possibility of the quantum term of the energy density becoming negative. 
However, when $B_{\rm cm} > \mean{B}_{\rm e}$, such as in a local maximum of the field, 
there is an amplification of the energy relative to the constant field case. We can put a 
bound on the difference between the mean field through a loop and the field at the 
center of mass for small loops (\ie small $T$),
\be
	|\mean{B} - B_{\rm cm}| \lesssim |B''(\rho_0)| T
\ee
where $|B''(\rho_0)| \ge B''(\rho)$ for all $\rho$ in the loop. This expression is proved the 
same way as determining the error in numerical integration using the midpoint rectangle 
rule.


If the field varies rapidly about some mean value 
on the Compton wavelength scale, the various contributions 
from local minima and local maxima are averaged out and the mean field approximation 
provided by the \ac{LCF} method becomes appropriate. A similar argument applies in the 
fermion case, where the important quantity is the mean magnetic field along the circumference 
of the loop. This quantity is also well served by a mean-field approximation when integrating 
over rapidly varying fields.

Another interesting feature of figure \ref{fig:EAscffresid} is that the \ac{LCF} 
approximation appears to describe narrower flux tubes better than wider ones, even when the 
spacing between the flux tubes is held constant. This effect is likely a result of the compact 
support given to the flux tube profiles. For narrow tubes, we are guaranteed to have many more 
center of mass points outside the flux tube than inside, giving a smaller energy contribution than 
for isolated flux tubes without compact support where the distinction between inside and outside is 
not as abrupt. This also explains why narrow, closely spaced tubes produce a lower energy than 
is predicted by the \ac{LCF} approximation.

This argument does not apply to the smooth isolated flux tubes given 
by equation (\ref{eqn:isoB}). 
For these flux tubes, the only region where there is a large discrepancy between $\mean{B}_{\rm e}$ and
$B_{\rm cm}$ is near the center of the flux tube. This 
is a global maximum of the field, and the only maximum of $|B''(\rho)|$. 
There are no regions where the average field in the loop ensemble is 
much stronger than the center of mass magnetic field. So, we expect an amplification of the 
energy near the flux tube relative to the constant field case. 
In the flux tubes with compact support, however, 
there is such a region just outside the flux tube. We can understand the surprisingly close agreement 
of these results to the \ac{LCF} approximation in our model 
in terms of competition between these 
regions of local minima and maxima of the field (see figure \ref{fig:EAscffvsconst}). 


\begin{figure}
	\centering
		\includegraphics[width=10cm]{images/EAscffvsconst.eps}
	\caption[A comparison between the action density in the \ac{LCF} approximation and \ac{WLN}]
	{The action density in \ac{WLN} and the \ac{LCF} approximation, scaled by $\rho$ so 
	an area on the figure is proportional to a volume. The approximation 
	is poor everywhere, however, when there are regions of local minima and local maxima of the field about 
	a mean value, the effective action approximately agrees between these methods. This is due to 
	a partial cancellation between regions where the estimate provided by the approximation is too 
	large (green) and other regions where the estimate is too small (red).}
	\label{fig:EAscffvsconst}
\end{figure}





Finally, I find that the quantum term remains small compared to the 
classical action for the range of parameters investigated. This is shown 
in figure \ref{fig:ActRatscff} where we plot the ratio of the \acs{ScQED}
term of the action to the classical action. The relative smallness of this correction 
is consistent with the predictions from homogeneous fields and the derivative 
expansion, as well as with previous studies on flux tube 
configurations~\cite{1999PhRvD..60j5019B}.
\begin{figure}
	\centering
		\includegraphics[width=10cm]{images/ActRatscff.eps}
	\caption[The 1-loop \acs{ScQED} quantum to classical ratio versus flux tube width]
	{The 1-loop \ac{ScQED} term divided by the classical term of the effective action as a function 
	of flux tube width, $\lambda/a$ for several values of the flux tube spacing, $a$.}
	\label{fig:ActRatscff}
\end{figure}

\subsection{Interaction Energies}

Using this model, we may also investigate the energy associated with interactions 
between the flux tubes. Since the flux tubes in our model exhibit compact support, 
the interaction energy is entirely due to a nonlocal interaction between nearby 
flux tubes. Thus, it contrasts with previous research which has investigated the 
interaction energies between flux tubes which have 
overlapping fields~\cite{Langfeld:2002vy}. In this 
case, there is a classical interaction energy ($\propto B_1 B_2$) as well as 
a quantum correction ($\propto (B_1 + B_2)^4 - B_1^4 - B_2^4$ in the weak-field limit). 
Even when these field overlap interactions are not present, 
there are also nonlocal energies in the 
vicinity of a flux tube due to the presence of other flux tubes. For example, 
the energy from nearby flux tubes can interact with a flux tube through the 
quantum diffusion of the magnetic field. Because of this 
phenomenon, we expect an interaction energy in the region of the central flux tubes 
due to the proximity of neighbouring flux tubes, even though no changes are made 
to the field profile or its derivatives in the region of interest. 
Since this interaction represents a force due to 
quantum fluctuations under the influence 
of external conditions, it is an example of a Casimir force. The Casimir force between 
two infinitely thin flux tubes in \ac{ScQED} has previously been found to be attractive~\cite{duru1993}. Our model 
can shed light specifically on this interaction, which is not predicted by 
local approximations such as the derivative expansion.

Consider a central flux tube with a width $\lambda = 0.5\lambda_e$. 
When $\lambda_{\rm min} = \lambda$, 
the magnetic field outside of the flux tube, $B_{\rm bg}$, is zero. Then, if the distance 
between flux tubes, $a$, is set very large, the energy density will be 
localized to the central flux tube and there will be no nonlocal interaction 
energy due to neighbouring flux tubes.
We define the interaction energy, $E_{\rm int}$, as the 
difference in energy within a distance $a/2$ of the central flux tube
between a configuration with a given value of $a$ 
and a configuration with $a = \infty$. In practice, we use $a = 10,000\lambda_e$ as 
a suitable stand-in for $a = \infty$:

\be
	\frac{E_{\rm int}}{L_z} = -\frac{\Gamma_{\rm scal}(a)}{L_z \mathcal{T}} 
	+\frac{\Gamma_{\rm scal}(a=1\times 10^4 \lambda_e)}{L_z \mathcal{T}}.
\ee
With this definition, the interaction energy is the energy associated with 
lowering the distance between flux rings from infinity. This the the analogue 
in our model of reducing the lattice spacing of the flux tubes.

One complication of this definition is that there is no clear distinction between 
energy density which `belongs' to the central flux tube and energy density which 
`belongs' to the neighbouring flux tubes. We continue to use our convention that the total energy 
for the central flux tube is determined by the integral over the non-local action density
in a region within a radius 
of $a/2$ of the flux tube. As $a$ is taken smaller and smaller, some energy from nearby 
flux tubes is included within this region, but also, some energy associated with the central 
flux tube is diffused out of the region. This ambiguity is unavoidable within this 
model. We can't numerically compute the energy over all of space and subtract off 
different contributions, because these energies are infinite.

\begin{figure}
	\centering
		\includegraphics[width=10cm]{images/interactionb.eps}
	\caption[The interaction energy versus the flux tube spacing]
	{The interaction energy per unit length of flux tube as a function 
	of the flux tube spacing, $a$. The energy density of a critical 
	strength magnetic field is 17 GeV/$\lambda_e^3$, so this 
	energy density is small in comparison to the classical magnetic field energies, 
	or the 1-loop corrections. However, the local interactions are constant in $a$.
	At $a<\lambda_e/2$, 
	the bump functions from neighbouring tubes overlap. This approximately corresponds
	to the critical magnetic field which destroys superconductivity.}
	\label{fig:interaction}
\end{figure}

The interaction energy is plotted in figure \ref{fig:interaction}. In this plot, the 
error bars are 1-sigma error bars that account for the correlations in the means 
computed for each group of worldlines,

\be
	\sigma_{E_{\rm int}} = \sqrt{\sigma_{E_a}^2 + \sigma_{E_{a=10000}}^2 
		- 2\operatorname{Cov}(E_a, E_{a=10000})},
\ee
where $\operatorname{Cov}(a,b)$ is the covariance between random variables $a$ and $b$.
Recall from figure \ref{fig:EAscffpeek} that there is a positive contribution 
to the effective action, and therefore a negative contribution to the energy 
from the region between flux tubes. As we reduce $a$, bringing the flux tubes 
closer together, two considerations become important. Firstly, we are increasing 
the average field strength meaning there tends to be more flux through the 
worldline loops which tends to give a negative contribution to the interaction 
energy. Secondly, we are reducing the spatial volume over which we integrate the 
energy since we only integrate $\rho$ from $0$ to $a/2$. This effect makes a positive 
contribution to the interaction energy since we include less and less of the 
region of negative energy density in our integral. 

In figure \ref{fig:interaction}, there appears to be a landscape with both 
positive and negative interaction energies at different values of $a$. 
These appear to be consistent with the interplay between positive and 
negative contributions described in the previous paragraph. This is consistent with 
the usual expectation of attractive Casimir forces~\cite{duru1993}. The dominant contribution 
for the positive energy values is caused by less of the negative energy region contributing 
as the domain assigned to the flux tube is reduced. However, the point at $a/\lambda_e = 2$ is 
negative ($-0.3\pm0.1$MeV/$\lambda_e$) indicating that an attractive interaction from 
nearby flux tubes is dominant. 
We are continuing to compute points at larger values of $a$. 

At $a/\lambda_e = 0.5$, 
the flux tubes are positioned right next to one another, and the negative 
contribution from the increase in the mean field appears to be 
slightly larger than the positive contribution from the loss of the 
region of negative energy density from the integral. Beyond
this, the flux tubes would overlap each other, which approximately corresponds 
to the critical background field which destroys superconductivity.

Based on the above explanation, it appears that the non-local 
interaction energy between magnetic fields has a strong dependence on 
the specific profile of the classical magnetic field that was used. 
This makes it difficult to predict if it will result in attractive 
or repulsive forces
in a more realistic model of a flux tube lattice. 

The energy density of a 
critical strength magnetic field is about $17{\rm GeV}/\lambda_e^3$, 
so the energy associated with this interaction is relatively small. 
However, there are no other interactions which affect the 
energy of moving the flux tubes closer together when they are 
separated by many coherence lengths. Here, the characteristic 
distance associated with the interaction, $\lambda_e$, is considerably 
larger than coherence length or London penetration depth, so the 
interactions between flux tubes through the order field are heavily 
suppressed. 


\section{Discussion and Conclusions}

In this chapter, I have developed a cylindrically symmetric magnetic field 
model which reproduces some of the 
features of a flux tube lattice: for a given central flux tube, there are 
nearby regions of large magnetic field which interact nonlocally, and 
the large flux tube size limit goes to a large uniform magnetic field
instead of to zero field. I have investigated the 1-loop effects from 
\ac{ScQED} in this model using the \ac{WLN} technique for various 
combinations of flux tube size, $\lambda$, and flux tube spacing, $a$.

In contrast to isolated flux tubes, I find that there are some regions where the 
\ac{WLN} results are greater than the \ac{LCF} approximation and other 
regions where they are less than the \ac{LCF} approximation. This can be understood 
by thinking of the difference from the \ac{LCF} approximation as a competition 
between the local counter term and the Wilson loop averages. 
For magnetic fields which vary on the Compton wavelength scale 
about some mean field strength, the \ac{LCF} approximation 
provides a poor approximation of the energy density, but may provide a 
good approximation to the total energy density of the field due to it 
being a good mean field theory approximation to the energy density. 
The appropriateness of the \ac{LCF} approximation in this case
can be understood as an approximate balance between regions where the 
field is a local maximum and the magnitude of the quantum correction to the 
action density is larger than in the constant field case, and regions where 
the field is a local minimum and the quantum corrections predict a smaller 
action density than the constant field case. This 
washing out of the field structure due to non-local effects has also 
been observed in \ac{WLN} studies of the vacuum polarization 
tensor~\cite{PhysRevD.84.065035}. 

There is a force between nearby magnetic flux tubes due 
to the quantum diffusion of the energy density. This interaction is 
non-local and is not predicted by the local derivative expansion.
It is an example of a Casimir force (\ie a force resulting from 
quantum vacuum fluctuations) and it is computed in a very similar 
way as the Casimir force between conducting bodies in the \ac{WLN} 
technique~\cite{Gies:2003cv}.
The size of the energy densities involved in this force are small even 
compared to the 1-loop corrections to the energy densities, which are
in turn small compared to the classical magnetic energy density. 

Although this interaction energy is small, 
the interactions between flux tubes in a neutron star due to the 
order field of the superconductor are suppressed because the distance between the 
tubes is considerably larger than the coherence length and 
London penetration depth. Therefore, this force is possibly important 
for the behaviour of flux tubes in neutron star crusts and interiors.
For example, in our lattice model, this force could contribute 
to a bunching of the worldlines, producing regions where flux tubes 
are separated by $\sim2\lambda_e$ and other regions which have 
no flux tubes. Consequently, this force may have important implications 
for neutron star physics. However, investigating these implications is 
outside the scope of this thesis.

The nature of this 
interaction energy is expected to depend on the model of the 
magnetic field profile for the reasons discussed in section 
\ref{sec:periodic_results}. It is therefore reasonable that 
forces of either sign, attractive or repulsive, may be possible 
depending on the particular landscape formed by the 
magnetic field and the particular definition used of the interaction 
energy. In a superconductor, the profiles of the magnetic 
flux tubes are determined by the minimization of the free 
energy for the interacting system formed by the magnetic 
and order parameter fields. Therefore, investigating this 
phenomenon using more realistic models is an interesting 
direction for future research. In particular, it would be interesting 
to determine if certain conditions allowed for a non-local interaction 
between magnetic flux tubes to be experimentally observable.

The conclusions from this chapter are applicable to \ac{ScQED}. 
However, we have also investigated the relationship between spinor 
and scalar \ac{QED} for isolated flux tubes where the \ac{WLN} technique 
can be applied to both cases. We find that both theories have the same 
qualitative behaviour, and agree within a factor of order unity quantitatively. The arguments 
and explanations given for the \ac{ScQED} results have strong parallels in the spinor 
\ac{QED} case. The spinor case can also be understood in terms of a competition between 
the Wilson loop averages and the local counter term. We therefore speculate that 
the results from this chapter will hold in the spinor case, at least qualitatively. However, 
addressing the fermion problem so that the spinor case can be studied explicitly for 
flux tube lattices would be valuable progress in this area of research.

